% GENERATED FILE. Edit canonical lessons, then run npm run generate:books.
% canonical-source-hash: fnv1a64:b99f6ead5ed8d352
% canonical-lessons: RU-C400-parol, RU-C400-bezopasnost, RU-C400-khaker, RU-C400-vzlamyvat, RU-C400-virus, RU-R400-review-393-400

\chapter{Staying safe online}
\label{ch:ru-staying-safe-online}
\hlchaptermodality{400}

\begin{chapteropening}
\textbf{By the end of this chapter:} \emph{I can talk about safety online: a password, a hacker, breaking into an account, and a virus.}

Complete the last lesson of chapter 400: I can talk about safety online: a password, a hacker, breaking into an account, and a virus.
\end{chapteropening}

\section[\texorpdfstring{\ru{пароль} (paról)}{пароль (paról)}]{\ru{пароль} (paról) — a password}

\label{lesson:RU-C400-parol}



\begin{quote}
Before the new one: say the Russian for a web browser, then the Russian for a search engine.
\end{quote}

\subsection*{You'll want to know: \ru{пароль}}
\textbf{\ru{пароль}} (\emph{paról}) — \textquotedblleft{}a password\textquotedblright{}.

\begin{quote}
\ru{Я} \ru{забыл} \ru{пароль} \ru{от} \ru{почты}.
\end{quote}

\begin{quote}
\emph{Ya zabýl paról ot póchty.}
\end{quote}

\begin{quote}
I have forgotten my email password.
\end{quote}

Masculine. \emph{\ru{Пароль} \ru{от}} + genitive — the password for.

\begin{cousinweb}
\textbf{\ru{пароль}} was borrowed from French \emph{parole} \textquotedblleft{}a word, a word of honour\textquotedblright{} (EtymDB), which is usually traced to Latin \emph{parabola} \textquotedblleft{}a comparison\textquotedblright{} (Lewis \& Short), later \textquotedblleft{}a saying\textquotedblright{}: first the word a sentry asks for.
\end{cousinweb}

\begin{grammarlens}[title={What you've built}]
One more word for telling, explaining and describing at B1.
\end{grammarlens}

\subsection*{Your turn}
\begin{itemize}
\raggedright
  \item \emph{Say it:} \emph{paról}
  \item \emph{Say it:} \emph{paról}, once more
  \item \emph{Say it:} a sentence of your own with \emph{paról}
  \item \emph{Recall:} say the Russian for a web browser, then the Russian for a search engine, then say \emph{paról} again
\end{itemize}

\begin{tcolorbox}[breakable,colback=teal!4,colframe=teal!35!black,title={Before you move on}]
What does \ru{пароль} mean? (\textquotedblleft{}A password\textquotedblright{}.) Say it once more.
\end{tcolorbox}
\section[\texorpdfstring{\ru{безопасность} (bezopásnost)}{безопасность (bezopásnost)}]{\ru{безопасность} (bezopásnost) — safety, security}

\label{lesson:RU-C400-bezopasnost}



\begin{quote}
Before the new one: say the Russian for a search engine, then the Russian for a password.
\end{quote}

\subsection*{You'll want to know: \ru{безопасность}}
\textbf{\ru{безопасность}} (\emph{bezopásnost}) — \textquotedblleft{}safety, security\textquotedblright{}.

\begin{quote}
\ru{Безопасность} \ru{детей} — \ru{главное}.
\end{quote}

\begin{quote}
\emph{Bezopásnost detéy — glávnoye.}
\end{quote}

\begin{quote}
Children’s safety comes first.
\end{quote}

Feminine. \emph{\ru{В} \ru{безопасности}} — safe. Built on \emph{\ru{опасный}}, dangerous.

\begin{cousinweb}
\textbf{\ru{безопасность}} is built from \emph{\ru{без}} \textquotedblleft{}without\textquotedblright{} and \emph{\ru{опасность}} \textquotedblleft{}danger\textquotedblright{}, from \emph{\ru{опасный}}, which is built on \emph{\ru{опасаться}} \textquotedblleft{}to be wary\textquotedblright{}, from \emph{\ru{пасти}} \textquotedblleft{}to guard\textquotedblright{}.
\end{cousinweb}

\begin{grammarlens}[title={What you've built}]
One more word for telling, explaining and describing at B1.
\end{grammarlens}

\subsection*{Your turn}
\begin{itemize}
\raggedright
  \item \emph{Say it:} \emph{bezopásnost}
  \item \emph{Say it:} \emph{bezopásnost}, once more
  \item \emph{Say it:} a sentence of your own with \emph{bezopásnost}
  \item \emph{Recall:} say the Russian for a search engine, then the Russian for a password, then say \emph{bezopásnost} again
\end{itemize}

\begin{tcolorbox}[breakable,colback=teal!4,colframe=teal!35!black,title={Before you move on}]
What does \ru{безопасность} mean? (\textquotedblleft{}Safety, security\textquotedblright{}.) Say it once more.
\end{tcolorbox}
\section[\texorpdfstring{\ru{хакер} (kháker)}{хакер (kháker)}]{\ru{хакер} (kháker) — a hacker}

\label{lesson:RU-C400-khaker}



\begin{quote}
Before the new one: say the Russian for a password, then the Russian for safety.
\end{quote}

\subsection*{You'll want to know: \ru{хакер}}
\textbf{\ru{хакер}} (\emph{kháker}) — \textquotedblleft{}a hacker\textquotedblright{}.

\begin{quote}
\ru{Хакер} \ru{украл} \ru{данные} \ru{клиентов}.
\end{quote}

\begin{quote}
\emph{Kháker ukrál dánnyye kliyéntov.}
\end{quote}

\begin{quote}
A hacker stole the customers’ data.
\end{quote}

Masculine. From English \emph{hacker}.

\begin{cousinweb}
\textbf{\ru{хакер}} was borrowed from English \emph{hacker} (EtymDB), keeping the English stress on the first syllable.
\end{cousinweb}

\begin{grammarlens}[title={What you've built}]
One more word for telling, explaining and describing at B1.
\end{grammarlens}

\subsection*{Your turn}
\begin{itemize}
\raggedright
  \item \emph{Say it:} \emph{kháker}
  \item \emph{Say it:} \emph{kháker}, once more
  \item \emph{Say it:} a sentence of your own with \emph{kháker}
  \item \emph{Recall:} say the Russian for a password, then the Russian for safety, then say \emph{kháker} again
\end{itemize}

\begin{tcolorbox}[breakable,colback=teal!4,colframe=teal!35!black,title={Before you move on}]
What does \ru{хакер} mean? (\textquotedblleft{}A hacker\textquotedblright{}.) Say it once more.
\end{tcolorbox}
\section[\texorpdfstring{\ru{взламывать} (vzlámyvat)}{взламывать (vzlámyvat)}]{\ru{взламывать} (vzlámyvat) — to break into, to hack}

\label{lesson:RU-C400-vzlamyvat}



\begin{quote}
Before the new one: say the Russian for safety, then the Russian for a hacker.
\end{quote}

\subsection*{You'll want to know: \ru{взламывать}}
\textbf{\ru{взламывать}} (\emph{vzlámyvat}) — \textquotedblleft{}to break into, to hack\textquotedblright{}.

\begin{quote}
\ru{Кто}-\ru{то} \ru{взломал} \ru{мою} \ru{почту}.
\end{quote}

\begin{quote}
\emph{Kto-to vzlomál moyú póchtu.}
\end{quote}

\begin{quote}
Someone has hacked my email.
\end{quote}

Imperfective; perfective \emph{\ru{взломать}}. Of doors and of accounts alike.

\begin{cousinweb}
\textbf{\ru{взламывать}} is built from \emph{\ru{вз}-} and \emph{\ru{ломать}} \textquotedblleft{}to break\textquotedblright{}, inherited from Proto-Slavic \emph{*lomiti}; the \emph{\ru{о}} turns to \emph{\ru{а}} in the imperfective.
\end{cousinweb}

\begin{grammarlens}[title={What you've built}]
One more word for telling, explaining and describing at B1.
\end{grammarlens}

\subsection*{Your turn}
\begin{itemize}
\raggedright
  \item \emph{Say it:} \emph{vzlámyvat}
  \item \emph{Say it:} \emph{vzlámyvat}, once more
  \item \emph{Say it:} a sentence of your own with \emph{vzlámyvat}
  \item \emph{Recall:} say the Russian for safety, then the Russian for a hacker, then say \emph{vzlámyvat} again
\end{itemize}

\begin{tcolorbox}[breakable,colback=teal!4,colframe=teal!35!black,title={Before you move on}]
What does \ru{взламывать} mean? (\textquotedblleft{}To break into, to hack\textquotedblright{}.) Say it once more.
\end{tcolorbox}
\section[\texorpdfstring{\ru{вирус} (vírus)}{вирус (vírus)}]{\ru{вирус} (vírus) — a virus}

\label{lesson:RU-C400-virus}



\begin{quote}
Before the new one: say the Russian for a hacker, then the Russian for to break into.
\end{quote}

\subsection*{You'll want to know: \ru{вирус}}
\textbf{\ru{вирус}} (\emph{vírus}) — \textquotedblleft{}a virus\textquotedblright{}.

\begin{quote}
\ru{На} \ru{компьютере} \ru{вирус}.
\end{quote}

\begin{quote}
\emph{Na kompyútere vírus.}
\end{quote}

\begin{quote}
There is a virus on the computer.
\end{quote}

Masculine. For computers and for illness alike.

\begin{cousinweb}
\textbf{\ru{вирус}} is usually traced to Latin \emph{vīrus} \textquotedblleft{}slime, poison\textquotedblright{} (Lewis \& Short); the computer sense came through English.
\end{cousinweb}

\begin{grammarlens}[title={What you've built}]
One more word for telling, explaining and describing at B1.
\end{grammarlens}

\subsection*{Your turn}
\begin{itemize}
\raggedright
  \item \emph{Say it:} \emph{vírus}
  \item \emph{Say it:} \emph{vírus}, once more
  \item \emph{Say it:} a sentence of your own with \emph{vírus}
  \item \emph{Recall:} say the Russian for a hacker, then the Russian for to break into, then say \emph{vírus} again
\end{itemize}

\begin{tcolorbox}[breakable,colback=teal!4,colframe=teal!35!black,title={Before you move on}]
What does \ru{вирус} mean? (\textquotedblleft{}A virus\textquotedblright{}.) Say it once more.
\end{tcolorbox}
\section[(dialogue)]{Words from chapters 393–400 — a review}

\label{lesson:RU-R400-review-393-400}



\begin{quote}
Say the Russian for an application and to submit.
\end{quote}

\subsection*{Your turn}
Say these aloud:

\begin{itemize}
\raggedright
  \item an application, to submit, a notary, a power of attorney, a will
  \item a debt, a loan, to lend, to give back, a per cent
  \item savings, to invest, a deposit, cost, a budget
  \item to go up in price, to get cheaper, inflation, an exchange rate, an exchange
  \item to pay for, a payment, free, paid, profitable
  \item rent, to rent, a mortgage, to insure, insurance
  \item a network, a website, a link, a web browser, a search engine
  \item a password, safety, a hacker, to break into, a virus
\end{itemize}

\begin{tcolorbox}[breakable,colback=teal!4,colframe=teal!35!black,title={Before you move on}]
An application? (\textbf{\ru{заявление}}.) To submit? (\textbf{\ru{подавать}}.) A notary? (\textbf{\ru{нотариус}}.) A power of attorney? (\textbf{\ru{доверенность}}.) A will? (\textbf{\ru{завещание}}.) A debt? (\textbf{\ru{долг}}.) A loan? (\textbf{\ru{кредит}}.) To lend? (\textbf{\ru{одалживать}}.) To give back? (\textbf{\ru{возвращать}}.) A per cent? (\textbf{\ru{процент}}.) Savings? (\textbf{\ru{сбережения}}.) To invest? (\textbf{\ru{вкладывать}}.) A deposit? (\textbf{\ru{вклад}}.) Cost? (\textbf{\ru{стоимость}}.) A budget? (\textbf{\ru{бюджет}}.) To go up in price? (\textbf{\ru{дорожать}}.) To get cheaper? (\textbf{\ru{дешеветь}}.) Inflation? (\textbf{\ru{инфляция}}.) An exchange rate? (\textbf{\ru{курс}}.) An exchange? (\textbf{\ru{обмен}}.) To pay for? (\textbf{\ru{оплачивать}}.) A payment? (\textbf{\ru{платёж}}.) Free? (\textbf{\ru{бесплатный}}.) Paid? (\textbf{\ru{платный}}.) Profitable? (\textbf{\ru{выгодный}}.) Rent? (\textbf{\ru{аренда}}.) To rent? (\textbf{\ru{арендовать}}.) A mortgage? (\textbf{\ru{ипотека}}.) To insure? (\textbf{\ru{страховать}}.) Insurance? (\textbf{\ru{страховка}}.) A network? (\textbf{\ru{сеть}}.) A website? (\textbf{\ru{сайт}}.) A link? (\textbf{\ru{ссылка}}.) A web browser? (\textbf{\ru{браузер}}.) A search engine? (\textbf{\ru{поисковик}}.) A password? (\textbf{\ru{пароль}}.) Safety? (\textbf{\ru{безопасность}}.) A hacker? (\textbf{\ru{хакер}}.) To break into? (\textbf{\ru{взламывать}}.) A virus? (\textbf{\ru{вирус}}.)
\end{tcolorbox}
